%======================================================================
\chapter{Span zero}\label{ch:span}
%======================================================================

The \emph{span} of a continuum, introduced by Lelek, measures how far apart two points can stay while each of them runs
over the whole continuum (or one runs over the part traced by the other). Snake-like continua have span zero, and the
main theorem of Hoehn and Oversteegen (Chapters~\ref{ch:snake}--\ref{ch:ladders}) is a converse for hereditarily
indecomposable continua. This chapter and the next prove:

\begin{theorem}[Theorem~C of Hoehn--Oversteegen]\label{thm:theoremC}
Every homogeneous indecomposable plane continuum has span zero.
\end{theorem}

\formal{\lean{HPC.Span.TheoremC}, \lean{HPC.Span.theoremC}.}

The proof follows Oversteegen and Tymchatyn, \emph{Plane strips and the span of continua (I)} (1982): Theorem~3.2 (the strip
theorem, Chapter~\ref{ch:strips}) $\Rightarrow$ Corollary~3.5 $\Rightarrow$ Theorem~3.6 $\Rightarrow$ Theorem~C. Several
steps of the printed proofs are unargued or, as printed, impossible; the formal proof supplies them, and proves the strip
theorem in a slightly weaker form that suffices.

\section{Span}

Let $X$ be a metric space with projections $\pr1,\pr2\colon X\times X\to X$, and $\Diag$ the diagonal. Only positivity of span matters,
so we use topological forms.

\begin{definition}\label{def:span}
A set $Z\subseteq X\times X$ is \emph{off the diagonal} if $Z\cap\Diag=\emptyset$. A set $A\subseteq X$ has \emph{positive surjective span} if there is a
continuum $Z\subseteq A\times A$ off the diagonal with $\pr1(Z)=\pr2(Z)=A$, and \emph{positive surjective semi-span} if there is one with $\pr1(Z)=A$.
$X$ has \emph{span zero} (\emph{semi-span zero}) if no subcontinuum of $X$ has positive surjective span (semi-span). Following
Hoehn--Oversteegen, $X$ has \emph{span zero in their sense} if every subcontinuum $Z\subseteq X\times X$ with $\pr1(Z)\subseteq\pr2(Z)$ meets the diagonal.
\end{definition}

\formal{\lean{HPC.Span.OffDiag}, \lean{HPC.Span.SurjSpanPos}, \lean{HPC.Span.SurjSemiSpanPos},
\lean{HPC.Span.SpanZero}, \lean{HPC.Span.SemiSpanZero}, \lean{HPC.Span.HOSpanZero}.}

\begin{proposition}\label{prop:span-basic}
\leavevmode
\begin{enumerate}[label=(\alph*)]
\item For compact $X$, Lelek's numerical definitions agree with these: $\sigma_0(X)>0$ (respectively $\sigma(X)>0$) iff some subcontinuum has
positive surjective semi-span (span).
\item Positive surjective span (semi-span) of $A$ passes to $f(A)$ for any continuous $f$ injective on $A$.
\item Semi-span zero is the same as span zero in the sense of Hoehn--Oversteegen; and that is equivalent to: for all maps $f,g\colon C\to X$
of a continuum with $f(C)\subseteq g(C)$ there is $c$ with $f(c)=g(c)$.
\end{enumerate}
\end{proposition}

\begin{proof}
(a) Given a connected $C\subseteq A\times A$ with $d\ge\alpha>0$ on $C$ and $\pr1(C)=A$, the closure $\cl C$ is a continuum off the diagonal with
$\pr1(\cl C)=\cl A$; conversely, on a continuum off the diagonal $d$ has a positive minimum. (b) Use $(f\times f)(Z)$. (c) Use the swap
$(a,b)\mapsto(b,a)$: if $Z$ is off the diagonal with $\pr1(Z)\subseteq\pr2(Z)=A$, its swap lies in $A\times A$ and has first projection $A$; and
conversely. For maps, take $Z=\set{(f(c),g(c))}$, and conversely $C=Z$ with the two projections.
\end{proof}

\formal{\lean{HPC.Span.otSemiSpanPos_iff}, \lean{HPC.Span.otSpanPos_iff}, \lean{HPC.Span.SurjSpanPos.image},
\lean{HPC.Span.SurjSemiSpanPos.image}, \lean{HPC.Span.semiSpanZero_iff_hoSpanZero}, \lean{HPC.Span.hoSpanZero_iff_maps}.}

\section{Confluence and Davis's theorem}

\begin{theorem}[Cook 1967, Theorem~4]\label{thm:cook}
Let $f$ be a continuous map of a compact metric continuum $X$ onto a hereditarily indecomposable continuum $S$. Then $f$ is
\emph{confluent}: for every subcontinuum $Q\subseteq S$ and every component $C$ of $f^{-1}(Q)$, $f(C)=Q$.
\end{theorem}

\begin{proof}
If $C=X$ then $f(C)=S\supseteq Q$. Otherwise let $q\in Q$ and suppose $q\notin f(C)$. Then $f^{-1}(q)$ is compact and misses $C$, so a closed
$\delta$-neighbourhood $N$ of $C$ misses $f^{-1}(q)$; we may take $N\ne X$. Let $K$ be the component of a point $x\in C$ in $N$. By boundary
bumping (\vthree{Theorem~1.10}) $K$ meets the frontier of $N$, so $K\not\subseteq C$. Then $f(K)$ meets $Q$ (at $f(x)$) and is not inside $Q$
(otherwise $K$ would be a connected subset of $f^{-1}(Q)$ larger than $C$). Two subcontinua of the hereditarily indecomposable $S$ that meet are
nested, so $Q\subseteq f(K)$, and $q=f(z)$ with $z\in K\subseteq N$; but $N$ misses $f^{-1}(q)$.
\end{proof}

\formal{\lean{HPC.Span.nested_of_hi}, \lean{HPC.Span.cook67_theorem4}.}

This uses one closed neighbourhood of $C$ instead of Cook's decreasing sequence $C_n$, whose existence Cook does not argue.

\begin{lemma}[Brouwer reduction]\label{lem:brouwer}
Let $Z$ be a continuum and $P$ a property of subcontinua of $Z$ that holds for $Z$ and passes to intersections of nonempty chains. Then $Z$
contains a subcontinuum $Z_0$ with $P$ that is minimal: no proper subcontinuum of $Z_0$ has $P$. Moreover, the intersection of a nonempty chain
of continua is a continuum, and continuous images commute with such intersections.
\end{lemma}

\begin{proof}
The intersection of a nonempty chain of continua is compact, nonempty by the finite intersection property, and connected as a directed
intersection of continua. Zorn's lemma gives $Z_0$. For a continuous $\varphi$ and a chain $(W_j)$ of compact sets, $\varphi(\bigcap W_j)=\bigcap\varphi(W_j)$:
if $y\in\bigcap\varphi(W_j)$, the sets $W_j\cap\varphi^{-1}(y)$ form a chain of nonempty compacta.
\end{proof}

\formal{\lean{HPC.Span.isContinuum_sInter_chain}, \lean{HPC.Span.brouwer_reduction}.}

\begin{theorem}[Davis's theorem for hereditarily indecomposable continua]\label{thm:davis}
If $A$ is hereditarily indecomposable and has positive surjective semi-span, some subcontinuum $B\subseteq A$ has positive surjective span.
\end{theorem}

\begin{proof}
Let $Z\subseteq A\times A$ be a continuum off the diagonal with $\pr1(Z)=A$. The property $P(W)$: ``$\pr2(W)\subseteq\pr1(W)$'' holds for $Z$ and passes to
intersections of chains (Lemma~\ref{lem:brouwer}); let $Z_0\subseteq Z$ be minimal with it, and $B=\pr1(Z_0)$, a hereditarily indecomposable continuum.
Suppose $Q=\pr2(Z_0)\ne B$. Take $w\in Z_0$ with $\pr1(w)\in Q$ and let $Z'$ be the component of $w$ in $Z_0\cap\pr1^{-1}(Q)$. By Cook's theorem (applied to
$\pr1\colon Z_0\to B$), $\pr1(Z')=Q$; so $\pr2(Z')\subseteq Q=\pr1(Z')$ and $Z'\ne Z_0$, against minimality. Hence $\pr1(Z_0)=\pr2(Z_0)=B$.
\end{proof}

\formal{\lean{HPC.Span.exists_surjSpanPos_of_surjSemiSpanPos}, \lean{HPC.Span.semiSpanZero_of_spanZero}.}

\section{Theorem 3.6 from Corollary 3.5}

\begin{definition}
\emph{Corollary~3.5} is the statement: if $M$ is an indecomposable homogeneous plane continuum, every proper subcontinuum of $M$ has zero surjective
semi-span (that is, does not have positive surjective semi-span).
\end{definition}

\formal{\lean{HPC.Span.OTCor35}.}

\begin{lemma}\label{lem:dense-proper}
A nondegenerate indecomposable continuum has, for every $\delta>0$, a $\delta$-dense proper subcontinuum.
\end{lemma}

\begin{proof}
The proper subcontinua through a point $x$ are directed by union (Theorem~\ref{thm:composants}(b)) and their union, the composant, is dense
(Theorem~\ref{thm:composants}(a)). Finitely many of them meet each ball of a finite $\delta/2$-net, and their union is one proper subcontinuum.
\end{proof}

\formal{\lean{HPC.Span.union_proper_of_indec}, \lean{HPC.Span.exists_dense_proper_subcontinuum}.}

\begin{theorem}[Oversteegen--Tymchatyn, Theorem~3.6, indecomposable case]\label{thm:ot36}
Let $X$ be a homogeneous, indecomposable, hereditarily indecomposable continuum none of whose proper subcontinua has positive surjective
semi-span. Then $X$ does not have positive surjective semi-span.
\end{theorem}

\begin{proof}
Suppose $Z\subseteq X\times X$ is a continuum off the diagonal with $\pr1(Z)=X$; write $f=\pr1$, $g=\pr2$ on $Z$, and let $\alpha>0$ be the minimum of $d(f,g)$ on $Z$.

\emph{Reduction.} We may assume that $g(Z)\ne X$, or that no proper subcontinuum of $Z$ is mapped onto $X$ by $g$: if $g(Z)=X$, take by
Lemma~\ref{lem:brouwer} a minimal $Z_0\subseteq Z$ with $g(Z_0)=X$; if $f(Z_0)=X$ replace $Z$ by $Z_0$, and otherwise by the swap of $Z_0$ (whose
first projection is onto and second is not).

\emph{Effros.} Let $\delta$ be the Effros bound (\vthree{Theorem~7.15}) for $\alpha$, and $K$ a $\delta$-dense proper subcontinuum (Lemma~\ref{lem:dense-proper};
$X$ is nondegenerate since $Z$ exists).

\emph{Cook.} By Theorem~\ref{thm:cook}, a component $C$ of $f^{-1}(K)$ has $f(C)=K$; $C\ne Z$ since $K\ne X$.

\emph{The copy.} Take $y\in g(C)$, $x\in K$ with $d(x,y)<\delta$, and an $\alpha$-homeomorphism $h$ with $h(x)=y$. The subcontinua $h(K)$ and $g(C)$ meet at
$y$, so they are nested; let $Z^*$ be the larger. It is proper: $h(K)\ne X$, and $g(C)\ne X$ because either $g(Z)\ne X$, or $g$ is irreducible onto $X$
and $C\subsetneq Z$.

\emph{The contradiction.} $W=\set{(h(f(c)),g(c)):c\in C}$ is a continuum with projections $h(K)$ and $g(C)$; it is off the diagonal, since
$h(f(c))=g(c)$ would give $d(f(c),g(c))=d(f(c),h(f(c)))<\alpha$. If $g(C)\subseteq h(K)=Z^*$, $W$ witnesses positive surjective semi-span of $Z^*$; otherwise its
swap does. This contradicts the hypothesis.
\end{proof}

\formal{\lean{HPC.Span.exists_reduced_witness}, \lean{HPC.Span.not_surjSemiSpanPos_univ}.}

\begin{corollary}\label{cor:C-from-35}
Corollary~3.5 implies Theorem~C.
\end{corollary}

\begin{proof}
Let $M$ be a homogeneous indecomposable plane continuum. It is hereditarily indecomposable (Theorem~\ref{thm:HI}). By Corollary~3.5 and
Theorem~\ref{thm:ot36} it has semi-span zero, which is span zero in the sense of Hoehn--Oversteegen
(Proposition~\ref{prop:span-basic}(c)).
\end{proof}

\formal{\lean{HPC.Span.ot82_theorem36_of_cor35}, \lean{HPC.Span.theoremC_of_cor35}.}

\section{Corollary 3.5 from the strip theorem}

The strip theorem is proved in the next chapter in the following form.

\begin{definition}[The strip theorem, strip form]\label{def:otthm32}
Let $Y$ be a continuum, $X\subseteq Y$ a subcontinuum, $x\in X$ and $y\in Y$. Let $(K_i)_{i\in I}$ be pairwise disjoint continua in the strip
$[0,1]\times\R$, each meeting both lines $x=0$ and $x=1$, and $h_i\colon Y\to K_i$ continuous maps with $h_i(x)$ on the line $x=0$ and
$\pr1(h_i(z))<\pr1(h_i(y))$ for all $z\in X$ (here $\pr1$ is the first coordinate). Then only countably many of the sets $h_i(X)$ have positive surjective
span.
\end{definition}

\formal{\lean{HPC.Span.OTThm32}.}

Oversteegen and Tymchatyn's Theorem~3.2 does not have the continua $K_i$; it only asks that the images $h_i(Y)$ be pairwise disjoint. Its printed
proof relies on two ``we may suppose'' steps that are not available together, and on three images one of which separates the other two, which is
not argued. The strip form is what Corollary~3.5 uses.

\begin{lemma}[Re-embedding]\label{lem:reembed}
\leavevmode
\begin{enumerate}[label=(E\arabic*)]
\item For a nonempty compact $K\subseteq\R^2$ there is a homeomorphism $\Phi$ of the plane with $\Phi(K)\subseteq\set{x\ge0}$ meeting the line $x=0$ in exactly one
point: if $(u_0,v_0)$ has least first coordinate in $K$, take $\Phi(u,v)=(u-u_0+(v-v_0)^2,v)$.
\item \emph{(Finger move.)} If $U$ is a connected open set and $p,q\in U$, there is a homeomorphism of the plane fixing every point off $U$ and
taking $p$ to $q$.
\item If $A\subseteq\set{x\ge0}$ is a continuum that does not separate the plane and meets $x=0$ in at most one point, then $\set{x>0}\setminus A$ is connected.
\end{enumerate}
\end{lemma}

\begin{proof}
(E2) If $B(c,r)\subseteq U$ and $\nrm{z-c}<r/2$, the map $T(w)=w+\beta(w)(z-c)$ with $\beta(w)=\max(0,1-\nrm{w-c}/r)$ is a homeomorphism (the perturbation is
Lipschitz with constant less than $\frac12$, so $T$ is injective and, by the contraction principle, onto), $T(c)=z$, and $T$ is the identity off
$B(c,r)$. The set of points reachable from $p$ by compositions of such maps is open and closed in $U$.

(E3) Let $F$ be the part of $\set{x>0}$ above, below, or to the right of a large box containing $A$; it is connected and misses $A$. Every component
$C$ of $\set{x>0}\setminus A$ meets $F$: otherwise $C$ is bounded with frontier in $A\cup J$, $J$ the segment of the line $x=0$ next to $A$, so $A\cup J$
separates the plane; but $A\cap J$ has at most one point and neither $A$ nor $J$ separates, against Janiszewski's theorem (\vfour{Theorem~2.7}).
\end{proof}

\formal{\lean{HPC.Span.exists_homeomorph_unique_leftmost}, \lean{HPC.Span.exists_homeomorph_ball_step},
\lean{HPC.Span.exists_homeomorph_push}, \lean{HPC.Span.isConnected_rightHalf_diff}.}

\begin{lemma}[Crossing continua]\label{lem:crossing}
\leavevmode
\begin{enumerate}[label=(\alph*)]
\item \emph{(The cut-wire argument.)} If $P$ is a continuum with points on both sides of the strip $S=[a,b]\times\R$, some component of $P\cap S$ meets
both lines $x=a$ and $x=b$.
\item If $M$ is a nondegenerate indecomposable plane continuum with points left of $x=a$ and right of $x=b$, there is an uncountable family of
pairwise disjoint subcontinua of $M\cap S$ each meeting both lines.
\end{enumerate}
\end{lemma}

\begin{proof}
(a) Otherwise, by \vthree{Theorem~1.9} applied in the compact set $P\cap S$, there is a relatively clopen $F_1$ containing $P\cap\set{x=a}$ and missing
$P\cap\set{x=b}$. With $F_2=(P\cap S)\setminus F_1$, the sets $(P\cap\set{x\le a})\cup F_1$ and $(P\cap\set{x\ge b})\cup F_2$ are disjoint, closed, nonempty and
cover $P$.

(b) Each composant $\kappa$ is dense, so it contains a point $p$ left of $x=a$, and $\cmp Mp=\kappa$ contains a point $q$ right of $x=b$; a proper subcontinuum
$P\ni p,q$ lies in $\kappa$, and (a) gives $K_\kappa\subseteq P\cap S$. Distinct composants are disjoint and uncountably many
(Theorem~\ref{thm:composants}).
\end{proof}

\formal{\lean{HPC.Span.exists_crossing_subcontinuum}, \lean{HPC.Span.exists_uncountable_crossing}.}

\begin{theorem}[Oversteegen--Tymchatyn, Corollary~3.5]\label{thm:ot35}
The strip theorem implies Corollary~3.5.
\end{theorem}

\begin{proof}
Let $M$ be an indecomposable homogeneous plane continuum and $A_0\ne M$ a subcontinuum with positive surjective semi-span. $M$ is hereditarily
indecomposable (Theorem~\ref{thm:HI}).

\emph{1. Davis.} By Theorem~\ref{thm:davis} some $A_1\subseteq A_0$ has positive surjective span.

\emph{2. Re-embed.} By (E1) we may assume $M\subseteq\set{x\ge0}$ with a unique point $x_0$ on $x=0$ (positivity of span and all hypotheses are carried
by $\Phi$). By homogeneity move $A_1$ to a continuum $A\ni x_0$; $A\ne M$ and $A$ has positive surjective span.

\emph{3. Finger move.} $A$ does not separate the plane (Theorem~\ref{thm:nosep}), so by (E3) $U=\set{x>0}\setminus A$ is open and connected. Take
$p\in M\setminus A$ and $q$ to the right of $A$, and by (E2) a homeomorphism $T$ fixing everything off $U$ with $T(p)=q$. Replacing $M$ by $T(M)$, we may
assume in addition that some point of $M$ lies to the right of $A$. (These printed ``without loss of generality'' steps are not a matter of
scaling; this is how the formal proof achieves them.)

\emph{4. Constants.} Let $s<m$ be the largest first coordinates on $A$ and on $M$, $\eta=(m-s)/8$ and $r=s+4\eta$. Let $\delta$ be the Effros bound for $\eta$,
and choose $\varepsilon\in(0,r)$ with $M\cap\set{x\le\varepsilon}\subseteq B(x_0,\delta)$.

\emph{5. The family.} Lemma~\ref{lem:crossing}(b) gives an uncountable family $\mathcal K$ of disjoint subcontinua crossing $[\varepsilon,r]\times\R$.

\emph{6. Copies of $A$.} For $K\in\mathcal K$ pick $x_K\in K$ on $x=\varepsilon$ and an $\eta$-homeomorphism $h_K$ with $h_K(x_0)=x_K$. $h_K(A)$ and $K$ meet, so
they are nested; $K\subseteq h_K(A)$ is impossible, since $K$ reaches $x=r$ while $h_K(A)\subseteq\set{x\le s+\eta}$. So $h_K(A)\subseteq K$.

\emph{7. The continuum $B$ and the point $y$.} The continua $B_K=h_K^{-1}(K)$ all contain $A$, so they are pairwise nested, and $B=\bigcap_KB_K$ is a
continuum containing $A$. The sets $B_K\cap\set{x\ge r-\eta}$ form a chain of nonempty compacta, so their intersection contains a point $y\in B$. For $z\in
A$: $\pr1h_K(z)<s+\eta<r-2\eta<\pr1h_K(y)$.

\emph{8.} Rescale the strip $[\varepsilon,r]\times\R$ to $[0,1]\times\R$ and apply the strip theorem with $Y=B$, $X=A$, $x=x_0$, $y$, the continua $K\in\mathcal K$ and the
maps $h_K|_B$. Every $h_K(A)$ has positive surjective span (Proposition~\ref{prop:span-basic}(b)), so $\mathcal K$ would be countable.
\end{proof}

\formal{\lean{HPC.Span.cor35_core}, \lean{HPC.Span.ot82_cor35_of_thm32}, \lean{HPC.Span.theoremC_of_thm32}.}

\begin{remark}
In the printed proof of Corollary~3.5, the normalization $\pr1(A)=[0,\frac12]$, the existence of $\mathcal K$, and the continuum $B$ (``it is easy to see'') are
unargued, the supposition $K\cap(\set\varepsilon\times\R)=\set{x_K}$ is unjustified (and unnecessary), and the inclusion ``$[0,\frac58]\subseteq\pr1(h_K(B))$'' is
impossible as printed; step~7 gives what is needed, one point $y$ for all $K$.
\end{remark}

\section*{Exercises}

\begin{exercise}
Show that an arc has span zero, and that a simple closed curve has positive surjective span.
\end{exercise}

\begin{exercise}
Show that a continuous image of a continuum with span zero need not have span zero, but that span zero is a topological invariant.
\end{exercise}
